% BOOK-FILE-ID: f-d0affc7bd30e
% 付録：本編の置き換えを行列で捉え直す

\section{漸化式と行列}\label{節：漸化式と行列}

本編では,さまざまな漸化式を等比型に帰着させてきた.
その多くは,\bookterm{matrix}{行列}を使うと1つの見方でまとめられる.

\paragraphbf{行列を公比のように見る}
等比型$a_{n+1}=ra_n$は,同じ数$r$を繰り返し掛ける漸化式であった.
では,記録しておく値が2つ以上あるときはどうだろうか.
値の組をベクトル$\boldsymbol v_n$にまとめ,更新を行列$M$で書けば,
\[
  \begin{array}{c|c}
    \text{等比型}&\text{行列}\\ \hline
    a_{n+1}=ra_n&\boldsymbol v_{n+1}=M\boldsymbol v_n\\
    a_n=r^{n-1}a_1&\boldsymbol v_n=M^{n-1}\boldsymbol v_1
  \end{array}
\]
と,等比型と全く同じ形になる.
行列$M$を公比のように見るのである.\footnote{「行列の等比数列」という言い方は直感を助けるための表現であり,正式な用語ではない.}

ただし,\,$M^{n-1}$は数の累乗のようには計算できない.
そこで,\sectionref*{節：行列の基本}で見たように,\,$M$を掛けてもただの定数倍になる向き,すなわち\bookterm{eigenvector}{固有ベクトル}を探す.
その向きの上では,行列を掛けることが本当に数を掛けることになるからである.

この節では,本編で「特性方程式」と呼んだ方程式を,次の2つに呼び分けて読んでいく.
\begin{itemize}
  \item \textbf{不動点の方程式}：\,$x=px+q$のように,変わらない値を求める方程式.
  \item \textbf{固有値の方程式}：\,$\det(\lambda I-M)=0$のように,更新の倍率を求める方程式.
\end{itemize}
本編の特性方程式がどちらにあたるかは,型によって異なる.
1階の漸化式,2階の漸化式,連立漸化式の順に確かめていこう.

\subsection{1階の漸化式を解く}\label{節：一つの値を更新する}

\subsubsection{特性方程式型：定数を成分に含める}
本編の例題$a_{n+1}=2a_n+1$では,毎回定数$1$が加わる.
その$1$も,\,$a_n$と一緒に記録してしまおう.
二つ目の成分を常に$1$に保てば,足し算もその成分の定数倍として書ける.
\[
  \begin{pmatrix}2&1\\0&1\end{pmatrix}
  \begin{pmatrix}a_n\\1\end{pmatrix}
  =\begin{pmatrix}2a_n+1\\1\end{pmatrix}
\]
下の成分は毎回$1$のままで,上の成分だけが元の規則で更新される.

一般の$a_{n+1}=pa_n+q$ $(p\ne0,1)$でも,定数$q$を同じ更新に含めよう.\footnote{$p=0$なら,\,$n\ge2$で$a_n=q$と直接分かる.\,$p=1$は後で扱う.}
状態を$\boldsymbol{x}_n=\begin{pmatrix}a_n\\1\end{pmatrix}$とすると,
\[
  \begin{aligned}
    A&=\begin{pmatrix}p&q\\0&1\end{pmatrix},\qquad
    \boldsymbol{x}_{n+1}=A\boldsymbol{x}_n,\\
    \boldsymbol{x}_n&=A^{n-1}\boldsymbol{x}_1
  \end{aligned}
\]
となる.

この$A$の固有ベクトルを探そう.
本編の不動点の方程式$\alpha=p\alpha+q$を使えば
\[
  A\begin{pmatrix}\alpha\\1\end{pmatrix}
  =\begin{pmatrix}p\alpha+q\\1\end{pmatrix}
  =\begin{pmatrix}\alpha\\1\end{pmatrix}
\]
となり,\,$(\alpha,1)$は固有値1の固有ベクトルである.一方,
\[
  A\begin{pmatrix}1\\0\end{pmatrix}
  =p\begin{pmatrix}1\\0\end{pmatrix}
\]
なので,\,$(1,0)$は固有値$p$の固有ベクトルである.
初期状態をこの2つの向きに分けると
\[
  \begin{pmatrix}a_1\\1\end{pmatrix}
  =\begin{pmatrix}\alpha\\1\end{pmatrix}
  +(a_1-\alpha)\begin{pmatrix}1\\0\end{pmatrix}
\]
であり,\,$n-1$回更新して上の成分を読めば
\[
  a_n=\alpha+(a_1-\alpha)p^{n-1}
\]
を得る.本編の「\bookterm{fixed}{不動点}」と「ずれ」が,倍率1と倍率$p$の2つの向きになったのだ.

ここで2つの方程式の役割を確かめておこう.
固有値は,固有値の方程式
\[
  \det(\lambda I-A)
  =(\lambda-p)(\lambda-1)=0
\]
から$\lambda=1,p$と求まる.これは倍率である.
一方,本編の$\alpha=p\alpha+q$は不動点の方程式であり,その解$\alpha$は倍率ではなく,固有値1の固有ベクトル$(\alpha,1)$の\textbf{成分の比}を表している.
特性方程式型の「特性方程式」は,不動点の方程式だったのである.

量を取り出す見方もできる.本編で$b_n=a_n-\alpha$と置いたのは
\[
  b_n=(1,-\alpha)\boldsymbol x_n,\qquad (1,-\alpha)A=p(1,-\alpha)
\]
より,\bookterm{left-eigen}{左固有ベクトル}$(1,-\alpha)$で,毎回$p$倍される量を取り出す操作であった.

最後に$p=1$の場合を見ておく.
\[
  A=\begin{pmatrix}1&q\\0&1\end{pmatrix},
  \qquad
  A^{n-1}=\begin{pmatrix}1&(n-1)q\\0&1\end{pmatrix}
\]
なので$a_n=a_1+(n-1)q$となり,等差型に戻る.
固有値は1が重なり,\,$q\ne0$なら固有ベクトルは$(1,0)$の向きしかない.
このように固有値が重なって固有ベクトルが足りないとき,一般項に$n$が現れる.これは\sectionref*{節：固有ベクトルが足りないとき}で詳しく見よう.

\BookImport{f-5bb007eab8c1}

\subsubsection{一次分数型：二成分の比を読む}\label{節：分数を二つの成分で追う}

特性方程式型では,下の成分を$1$に固定していた.
では,その固定をやめて,2つの成分を自由に動かしたらどうなるだろうか.
\sectionref{節：一次分数型}の
\[
  a_{n+1}=\frac{2a_n}{1-a_n}
\]
で試してみよう.各項で$a_n\ne1$とし,\,$a_1\ne-1$とする.
1つの値を分子と分母に分けて持ち,\,$a_n=\frac{x_n}{y_n}$と書く.
$(x_1,y_1)=(a_1,1)$から
\[
  \begin{pmatrix}x_{n+1}\\y_{n+1}\end{pmatrix}
  =M\begin{pmatrix}x_n\\y_n\end{pmatrix},
  \qquad
  M=\begin{pmatrix}2&0\\-1&1\end{pmatrix}
\]
で更新すると,
\[
  \frac{x_{n+1}}{y_{n+1}}
  =\frac{2x_n}{y_n-x_n}
  =\frac{2a_n}{1-a_n}
\]
となり,元の更新に戻る.
分数の漸化式が,2つの成分の行列による更新とその比に分かれたのだ.

本編で用いた$b_n=\frac{a_n}{a_n+1}$は,この二成分では$b_n=\frac{x_n}{x_n+y_n}$である.
分子は1回ごとに$2$倍されるが,分母は
\[
  x_{n+1}+y_{n+1}=2x_n+(y_n-x_n)=x_n+y_n
\]
と変わらない.従って,\,$b_{n+1}=2b_n$となる.
$x_1+y_1=a_1+1\ne0$なので,この比の分母はずっと零でない.

\bookterm{fixed}{不動点}$0,-1$を二成分で表すと,
\[
  M\begin{pmatrix}0\\1\end{pmatrix}
  =\begin{pmatrix}0\\1\end{pmatrix},
  \qquad
  M\begin{pmatrix}-1\\1\end{pmatrix}
  =2\begin{pmatrix}-1\\1\end{pmatrix}
\]
となる.不動点は固有ベクトルの成分比であり,対応する固有値は$1,2$である.
行列で$(x_n,y_n)$を動かすのは一次変換であり,その成分の比$\frac{x_n}{y_n}$だけを読んだものが一次分数変換である.
比を読むと長さの情報は消えるので,固有ベクトルからは向き,すなわち成分の比だけが残る.これが不動点である.
一方,その向きが何倍されるかを表すのが固有値であり,2つの固有値の比$2:1$が,本編の差の比の公比$2$として現れていた.

\paragraphbf{一般の一次分数型へ}
一般の
\[
  a_{n+1}=\frac{pa_n+q}{ra_n+s}
\]
でも同じである.\,$r\ne0$, $ps-qr\ne0$を仮定し,各項で$ra_n+s\ne0$とする.
$(x_1,y_1)=(a_1,1)$とし,
\[
  \begin{pmatrix}x_{n+1}\\y_{n+1}\end{pmatrix}
  =M\begin{pmatrix}x_n\\y_n\end{pmatrix},
  \qquad
  M=\begin{pmatrix}p&q\\r&s\end{pmatrix}
\]
で更新すれば,\,$a_n=\frac{x_n}{y_n}$が元の漸化式を満たす.
$y_{n+1}=(ra_n+s)y_n$なので,元の漸化式の分母が零にならない限り,\,$y_n\ne0$が保たれる.
同じ零でない数で二成分を倍しても比は変わらない.この成分比を通して行列の作用を見る考え方を,射影変換という.

不動点$\alpha$では
\[
  M\begin{pmatrix}\alpha\\1\end{pmatrix}
  =(r\alpha+s)\begin{pmatrix}\alpha\\1\end{pmatrix}
\]
となる.ここでも,不動点$\alpha$そのものは固有値ではなく,固有ベクトルの成分比である.
異なる2つの不動点$\alpha,\beta$があり,対応する倍率を$\lambda_\alpha=r\alpha+s,\ \lambda_\beta=r\beta+s$とする.
状態を
\[
  \begin{pmatrix}x_n\\y_n\end{pmatrix}
  =c_n\begin{pmatrix}\alpha\\1\end{pmatrix}
  +d_n\begin{pmatrix}\beta\\1\end{pmatrix}
\]
と分けると,\,$c_n,d_n$はそれぞれ毎回$\lambda_\alpha,\lambda_\beta$倍される.
$a_1\ne\beta$とすれば,\,$M$は可逆で$\lambda_\alpha\ne0$だから$c_n\ne0$が保たれる.\footnote{$a_1=\beta$なら,\,$a_n=\beta$という定数列になる.}
成分を比較すると$x_n-\alpha y_n=(\beta-\alpha)d_n$,\,$x_n-\beta y_n=(\alpha-\beta)c_n$なので,
\[
  \frac{a_n-\alpha}{a_n-\beta}=-\frac{d_n}{c_n},
  \qquad
  \frac{a_{n+1}-\alpha}{a_{n+1}-\beta}
  =\frac{\lambda_\beta}{\lambda_\alpha}
  \frac{a_n-\alpha}{a_n-\beta}
\]
を得る.
本編の不動点の方程式$rx^2+(s-p)x-q=0$の解と係数の関係より$r(\alpha+\beta)=p-s$なので,
\[
  \lambda_\beta=p-r\alpha,\qquad
  \lambda_\alpha=p-r\beta
\]
となり,公比$\frac{\lambda_\beta}{\lambda_\alpha}=\frac{p-r\alpha}{p-r\beta}$は本編と一致する.
差の比の公比は,\textbf{2つの固有値の比}だったのである.

比を読むことの代償も確かめておこう.
比を取ると,2つの成分の大きさの情報は失われる.それでも困らないのは,一次分数型が比だけで閉じた漸化式だからである.
また,不動点が1つに重なるときは,固有ベクトルが1つしかなく,差の比は作れない.
本編でこの場合に逆数を取って等差型にしたのは,\sectionref*{節：固有ベクトルが足りないとき}と同じ状況なのである.

\subsection{2階の漸化式を解く}\label{節：二つの値を更新する}

\subsubsection{直前の2項を状態にする}
2階の漸化式では,次の項を決めるのに直前の2項が要る.そこで,その2項を状態として記録する.
例えば
\[
  x_{n+1}=x_n+y_n,
  \qquad
  y_{n+1}=x_n
\]
という2つの規則は
\[
  \begin{pmatrix}x_{n+1}\\y_{n+1}\end{pmatrix}
  =
  \begin{pmatrix}1&1\\1&0\end{pmatrix}
  \begin{pmatrix}x_n\\y_n\end{pmatrix}
\]
と1つの行列で表せる.
初めを$(x_0,y_0)=(1,0)$とすると
\[
  (1,0),\ (1,1),\ (2,1),\ (3,2),\ldots
\]
と進み,\,$x_n=F_{n+1},\ y_n=F_n$となる.
ここで$F_0=0,\ F_1=1$はフィボナッチ数列である.
1歩前の$x$を$y$に残すので,\,$x$だけを見れば$x_{n+2}=x_{n+1}+x_n$という2階の漸化式になる.

\sectionref{節：隣接3項型}も同じである.
\[
  a_{n+2}=sa_{n+1}+ta_n
\]
に対して
\[
  \boldsymbol{v}_n=
  \begin{pmatrix}a_{n+1}\\a_n\end{pmatrix},
  \qquad
  B=\begin{pmatrix}s&t\\1&0\end{pmatrix}
\]
とすると,
\[
  B\boldsymbol{v}_n
  =\begin{pmatrix}sa_{n+1}+ta_n\\a_{n+1}\end{pmatrix}
  =\boldsymbol{v}_{n+1}
\]
となる.二行目が直前の項を残すので,1階の更新$\boldsymbol{v}_{n+1}=B\boldsymbol{v}_n$になった.
本編の形$a_{n+2}+pa_{n+1}+qa_n=0$なら,\,$s=-p,\ t=-q$である.

$B$の固有ベクトルを$(x,y)$とすると,二行目から$x=\lambda y$となる.
$y=1$にそろえると固有ベクトルは$\begin{pmatrix}\lambda\\1\end{pmatrix}$であり,一行目から
\[
  \lambda^2-s\lambda-t=0
\]
を得る.これが$B$の固有値の方程式であり,本編の特性方程式$x^2-sx-t=0$(本編の符号では$x^2+px+q=0$)と同じ方程式である.
\textbf{隣接3項型の「特性方程式」は,固有値の方程式だった}のだ.
特性方程式型では不動点の方程式であったのと対照的である.

この2つは,実はつながっている.
$a_{n+1}=pa_n+q$の添字を1つずらして引くと,定数$q$が消えて
\[
  a_{n+2}-(p+1)a_{n+1}+pa_n=0
\]
という隣接3項型になる.\footnote{本編\sectionref*{節：特性方程式型}の別解で用いた変形である.}
この特性方程式は$(x-1)(x-p)=0$であり,特性方程式型の行列$A$の固有値の方程式と一致する.
添字をずらして定数を消すと,不動点の方程式から固有値の方程式へ移るのである.

2つの異なる\bookterm{root}{特性根}$\lambda_1,\lambda_2$があれば,固有ベクトル$(\lambda_1,1),(\lambda_2,1)$は同じ直線上にない.
初めの状態をこの2つの向きに分けておけば,\,$k$回の更新でそれぞれの部分が$\lambda_1^k,\lambda_2^k$倍される.
こうして一般項に特性根の冪が現れる.
フィボナッチ数列では$x^2-x-1=0$より
\[
  \varphi=\frac{1+\sqrt5}{2},
  \qquad
  \psi=\frac{1-\sqrt5}{2}
\]
であり,初期値$F_0=0,\ F_1=1$に合わせて2つの向きの割合を決めると
\[
  F_n=\frac{\varphi^n-\psi^n}{\sqrt5}
\]
を得る.\,$|\psi|<1<\varphi$だから,\,$\psi^n$は符号を交互に変えながら減衰し,\,$\varphi^n$が増大を支配する.
特性根は,数列の長期的な動きも表しているのである.

\subsubsection{量を取り出すか,状態を分けるか}
本編で作った$a_{n+1}-\lambda_1a_n$という量も,同じ行列の中に見つかる.
解と係数の関係$\lambda_1+\lambda_2=s$,\,$\lambda_1\lambda_2=-t$を使うと
\[
  \begin{aligned}
    (1,-\lambda_1)B
    &=(s-\lambda_1,t)\\
    &=(\lambda_2,-\lambda_1\lambda_2)
    =\lambda_2(1,-\lambda_1)
  \end{aligned}
\]
なので,
\[
  b_n=(1,-\lambda_1)\boldsymbol{v}_n=a_{n+1}-\lambda_1a_n
\]
は$b_{n+1}=\lambda_2b_n$を満たす.
本編の$a_{n+2}=5a_{n+1}-6a_n$で作った$a_{n+1}-2a_n$は,\,$(1,-2)$で取り出す$3$倍の量だった.
$a_{n+1}-\alpha a_n$を作る操作は,\bookterm{left-eigen}{左固有ベクトル}で\textbf{一定倍率で変わる量を取り出す}操作だったのである.

これに対し,初めの状態を固有ベクトルの向きに分けるのは,\textbf{状態そのものを分ける}操作である.
同じ更新を,量の側から見るか,状態の側から見るかの違いにすぎない.
ただし,取り出した量$b_n$は1つの数なので,それだけでは$a_n$に戻れない.
本編で$a_1$を使って$a_n$を求め直したり,もう一方の根で作った量と連立させたりしたのはこのためである.

\subsubsection{非実数の特性根と回転}\label{節：非実数の特性根と回転}
特性根が実数でないときはどうなるだろうか.具体例で見よう.
\[
  a_1=a_2=1,\qquad a_{n+2}=a_{n+1}-a_n
\]
を計算すると
\[
  1,\ 1,\ 0,\ -1,\ -1,\ 0,\ 1,\ 1,\ \ldots
\]
と,周期6で同じ値を繰り返す.
特性方程式$x^2-x+1=0$の解は
\[
  \omega=\frac{1+\sqrt3i}{2}=\cos60^\circ+i\sin60^\circ
\]
とその共役$\overline\omega$である.\,$60^\circ$を6回足すと1周するので,周期6はこの回転から来ていそうである.

ところが,状態$(a_{n+1},a_n)$を座標平面に描くと
\[
  (1,1),\ (0,1),\ (-1,0),\ (-1,-1),\ (0,-1),\ (1,0)
\]
となり,原点からの距離は$\sqrt2$と$1$が混じっている.点は円ではなく,楕円
\[
  a_{n+1}^2-a_{n+1}a_n+a_n^2=1
\]
の上を回っているのだ.\footnote{6点すべてで左辺が1になることを確かめよう.この値が更新で変わらないことは,次の$u_n^2+w_n^2$から分かる.}
元の座標のままでは,回転はゆがんで見えるのである.

そこで,前項の「量を取り出す」方法を複素数の根で使ってみよう.
$z_n=a_{n+1}-\overline\omega a_n$と置くと,前項と同じ計算で$z_{n+1}=\omega z_n$となる.
$\overline\omega=\frac12-\frac{\sqrt3}{2}i$なので,\,$z_n$の実部と虚部は
\[
  u_n=a_{n+1}-\frac12a_n,\qquad w_n=\frac{\sqrt3}{2}a_n
\]
である.\,$\omega$を掛けることは,\sectionref*{節：回転と複素数の掛け算}で見たように$60^\circ$の回転なので,
\[
  \begin{pmatrix}u_{n+1}\\w_{n+1}\end{pmatrix}
  =R(60^\circ)\begin{pmatrix}u_n\\w_n\end{pmatrix}
\]
となる.実際,\,$u_{n+1}=\frac12u_n-\frac{\sqrt3}{2}w_n$,\,$w_{n+1}=\frac{\sqrt3}{2}u_n+\frac12w_n$を直接確かめられる.
座標を$(u_n,w_n)$に選び直せば,点はちょうど円の上を$60^\circ$ずつ回るのだ.
$u_n^2+w_n^2=a_{n+1}^2-a_{n+1}a_n+a_n^2$なので,この円を元の座標に戻したものが先ほどの楕円である.

一般に,特性根が$r(\cos\theta\pm i\sin\theta)$なら,同じように実数の座標を選ぶと,更新は角$\theta$の回転と$r$倍の拡大・縮小になる.
$r=1$なら点は楕円の上を回り,\,$r\ne1$なら渦を巻いて外へ広がるか中心へ縮む.
\sectionref*{節：特性根が虚数の場合}で一般項に$r^n\cos n\theta,\ r^n\sin n\theta$が現れたのは,このためである.

特に$r=1$のとき,特性方程式は$x^2-2\cos\theta\,x+1=0$なので,漸化式は
\[
  a_{n+2}=2\cos\theta\,a_{n+1}-a_n
\]
となる.余弦の加法定理より,\,$a_n=\cos(n\theta)$はこれを満たす.
この漸化式は,\textbf{一定の角$\theta$を足していく}更新なのである.
先ほどの例は$\theta=60^\circ$の場合にあたる.

\subsubsection{重解：固有ベクトルが足りないとき}\label{節：固有ベクトルが足りないとき}
特性方程式が重解$\alpha\ne0$を持つと,固有ベクトルは$(\alpha,1)$の向きしかなく,2つの向きに分けられない.
このとき一般項に$n$が現れることは,本編\sectionref*{節：付加項がない場合}で見た.
本編では,\,$a_{n+1}-\alpha a_n$が等比型になることを使い,\,$c_n=\frac{a_n}{\alpha^{n-1}}$と倍率を取り除くと$c_n$の差が一定になった.
つまり,固有ベクトルに分けきれない部分は,\textbf{等比的な倍率を取り除いて差分を追う}ことで求められるのである.
この方法を3重,4重の重解や付加項がある場合へ広げることは,\sectionref*{節：差分で線形漸化式を解く}で扱う.

では,固有値が重なれば必ず$n$が現れるのだろうか.
連立型
\[
  x_{n+1}=2x_n,\qquad y_{n+1}=2y_n
\]
の更新行列$\begin{pmatrix}2&0\\0&2\end{pmatrix}$は,固有値2を重ねて持つ.
しかし$x_n=2^{n-1}x_1,\ y_n=2^{n-1}y_1$であり,\,$n$はどこにも現れない.
この行列ではどの向きのベクトルも2倍されるので,固有ベクトルが2方向とも取れるからである.
$n$が現れるのは,固有値が重なり,\textbf{しかも固有ベクトルが足りない}ときなのだ.
隣接3項型の行列$B$では固有ベクトルが$(\lambda,1)$の向きに限られるので,重解なら必ず足りなくなる.
特性方程式型の$p=1$,すなわち等差型も,同じ理由で$n$が現れた例である.

\subsection{連立漸化式を解く}\label{節：連立漸化式を解く}

\subsubsection{重みは左固有ベクトル}
\sectionref{節：連立型}では,\,$x_{n+1}=2x_n+y_n,\ y_{n+1}=x_n+2y_n$の和と差$u_n=x_n+y_n$,\,$v_n=x_n-y_n$が
\[
  u_{n+1}=3u_n,\qquad v_{n+1}=v_n
\]
を満たすことを使った.\sectionref*{節：行列の基本}の\bookterm{left-eigen}{左固有ベクトル}$(1,1)$は和を,\,$(1,-1)$は差を取り出していたのである.

一般に,連立漸化式を
\BookMathLayout{\[
  \boldsymbol{z}_{n+1}=A\boldsymbol{z}_n,\qquad
  \boldsymbol{z}_n=\begin{pmatrix}a_n\\b_n\end{pmatrix},\qquad
  A=\begin{pmatrix}p&q\\r&s\end{pmatrix}
\]}{\[
  \begin{aligned}
    &\boldsymbol{z}_{n+1}=A\boldsymbol{z}_n,\\
    &\boldsymbol{z}_n=\begin{pmatrix}a_n\\b_n\end{pmatrix},\quad
    A=\begin{pmatrix}p&q\\r&s\end{pmatrix}
  \end{aligned}
\]}
と書き,重み付きの和$c_n=\alpha a_n+\beta b_n$を考える.これは$c_n=(\alpha,\beta)\boldsymbol z_n$と書ける.
$(\alpha,\beta)A=\lambda(\alpha,\beta)$なら,\,$c_{n+1}=\lambda c_n$となる.
本編の連立型で$c_n$が等比数列になるように$\alpha,\beta$を選ぶ計算は,左固有ベクトルを求めていたのである.
そのような重みが存在する条件は固有値の方程式
\[
  \lambda^2-(p+s)\lambda+(ps-qr)=0
\]
であり,本編で現れた特性方程式と一致する.
連立型の「特性方程式」も,隣接3項型と同じく固有値の方程式であった.
異なる2つの固有値があれば,取り出した2つの量から$a_n,b_n$を求め直せる.

\subsubsection{座標軸を回すことと,点を回すこと}
先ほどの例で$(1,1)$と$(1,-1)$の向きを新しい座標軸にすると,座標軸を$45^\circ$回して見ることになる.
\BookFigSystem
回転行列
\[
  Q=\frac{1}{\sqrt2}
  \begin{pmatrix}
    1&-1\\
    1&1
  \end{pmatrix}
\]
で新しい座標$\boldsymbol w_n=Q^{\mathsf T}\boldsymbol z_n$を作ると,\,$Q^{\mathsf T}Q=QQ^{\mathsf T}=I$より$\boldsymbol z_n=Q\boldsymbol w_n$と戻せて,
\[
  \boldsymbol w_{n+1}
  =Q^{\mathsf T}AQ\boldsymbol w_n,
  \qquad
  Q^{\mathsf T}AQ=
  \begin{pmatrix}3&0\\0&1\end{pmatrix}
\]
となる.新しい座標では,2つの等比型に分かれるのだ.

ここで回しているのは\textbf{座標軸}であり,更新行列$A$が点を回しているわけではない.
$A$は,\,$(1,1)$の向きを3倍に引き伸ばし,\,$(1,-1)$の向きはそのままにする変形である.
これに対し,\sectionref*{節：非実数の特性根と回転}では\textbf{更新そのものが点を回していた}.
同じ「回転」という言葉でも,見る座標を変える話と,点が動く話は区別しておこう.

さらに,座標軸を回すだけで済んだのは,この例の2つの固有ベクトルがたまたま直交していたからである.\footnote{$A$が対称,すなわち$q=r$のときは,異なる固有値の固有ベクトルが必ず直交することが知られている.}
一般には固有ベクトルは直交せず,新しい座標は斜めの座標になる.次の例で確かめよう.

\subsubsection{1つの行列で見渡す}\label{節：1つの行列で見渡す}
最後に,1つの行列
\[
  M=\begin{pmatrix}1&2\\1&0\end{pmatrix}
\]
から,連立型・隣接3項型・一次分数型を同時に読んでみよう.
\begin{itemize}
  \item \textbf{連立型}として読めば,
    \[
      x_{n+1}=x_n+2y_n,\qquad y_{n+1}=x_n
    \]
    である.
  \item 2つ目の式を1つ目に代入すれば,
    \[
      x_{n+2}=x_{n+1}+2x_n
    \]
    という\textbf{隣接3項型}になる.
  \item 比$b_n=\frac{x_n}{y_n}$を読めば,
    \[
      b_{n+1}=\frac{x_n+2y_n}{x_n}=\frac{b_n+2}{b_n}
    \]
    という\textbf{一次分数型}になる.\footnote{$y_n=x_{n-1}$なので,\,$b_n$は隣接3項型の隣り合う項の比である.本編\sectionref*{節：隣接3項型}で$\frac{a_{n+1}}{a_n}$と置いて一次分数型を作った別解は,この読み方にあたる.}
\end{itemize}
1つの更新を,どの成分を見るか,比を見るかで読み分けているだけなのだ.

$M$の固有値の方程式は$\lambda^2-\lambda-2=(\lambda-2)(\lambda+1)=0$であり,
\[
  M\begin{pmatrix}2\\1\end{pmatrix}=2\begin{pmatrix}2\\1\end{pmatrix},
  \qquad
  M\begin{pmatrix}-1\\1\end{pmatrix}=-\begin{pmatrix}-1\\1\end{pmatrix}
\]
となる.
一方,一次分数型$b_{n+1}=\frac{b_n+2}{b_n}$の不動点の方程式$x=\frac{x+2}{x}$,すなわち$x^2-x-2=0$の解も$2,-1$である.
固有値と不動点が同じ値になったが,役割は異なる.
固有値$2,-1$は倍率であり,不動点$2,-1$は固有ベクトル$(2,1),(-1,1)$の成分比である.
この行列では固有ベクトルが$(\lambda,1)$の形をしているので,成分比がたまたま固有値と等しくなったにすぎない.
実際,差の比の公比は固有値の比$-\frac12$であり,
\[
  \frac{b_{n+1}-2}{b_{n+1}+1}=-\frac12\cdot\frac{b_n-2}{b_n+1}
\]
となる.\footnote{例えば$(x_1,y_1)=(1,1)$なら$b_1=1,\ b_2=3$で,\,$\frac{1-2}{1+1}=-\frac12$,\,$\frac{3-2}{3+1}=\frac14$である.}

固有ベクトル$(2,1)$と$(-1,1)$は直交しない.
それぞれの向きの成分を取り出す左固有ベクトルは
\[
  (1,1)M=2(1,1),\qquad (-1,2)M=-(-1,2)
\]
より$(1,1)$と$(-1,2)$であり,\,$\boldsymbol z=\begin{pmatrix}x\\y\end{pmatrix}$は
\[
  \boldsymbol z=\frac{x+y}{3}\begin{pmatrix}2\\1\end{pmatrix}+\frac{-x+2y}{3}\begin{pmatrix}-1\\1\end{pmatrix}
\]
と斜めの座標に分けられる.\footnote{左固有ベクトル$(1,1)$は,もう一方の固有ベクトル$(-1,1)$と直交している.$(-1,2)$と$(2,1)$も同様である.斜めの座標の成分は,このように「もう一方の向きに沿って測る」ことで取り出す.}
$(x_1,y_1)=(1,0)$なら,それぞれ$2^{n-1}$倍,\,$(-1)^{n-1}$倍して戻すと
\[
  x_n=\frac{2^n+(-1)^{n-1}}{3}
\]
を得る.\,$x_1=1,\ x_2=1,\ x_3=3$と,漸化式から計算した値に一致する.
座標軸を回すだけでは,この分解は作れないのである.

\BookImport{f-947cf9e311ee}

\subsection{一般の隣接多項間漸化式へ}\label{節：一般の隣接多項間漸化式へ}
記録する値を増やせば,同じ考え方がそのまま使える.
\sectionref{節：隣接4項型}の$a_{n+3}+pa_{n+2}+qa_{n+1}+ra_n=0$なら,直前の3項を状態にして
\[
  \begin{pmatrix}a_{n+3}\\a_{n+2}\\a_{n+1}\end{pmatrix}
  =\begin{pmatrix}-p&-q&-r\\1&0&0\\0&1&0\end{pmatrix}
  \begin{pmatrix}a_{n+2}\\a_{n+1}\\a_n\end{pmatrix}
\]
と書ける.下の2行は,直前の項を1つずつ送るだけである.
2階のときと同じく,固有ベクトルを$(\lambda^2,\lambda,1)$の形で探して一行目に代入すると
\[
  \lambda^3+p\lambda^2+q\lambda+r=0
\]
を得る.これは本編の隣接4項型の特性方程式$x^3+px^2+qx+r=0$と同じ方程式である.
三行三列の行列では行列式の計算が増えるが,ここでは使わない.\,$a_n=\lambda^n$を代入するだけで,同じ方程式が得られるからである.

一般に,直前の$k$項から次の項が決まる漸化式でも,
\begin{itemize}
  \item 直前の$k$項を状態にすれば,1階の更新$\boldsymbol v_{n+1}=M\boldsymbol v_n$になる.
  \item 固有ベクトルは$(\lambda^{k-1},\ldots,\lambda,1)$の形であり,固有値の方程式は$a_n=\lambda^n$を代入して得る特性方程式と一致する.
  \item 異なる特性根がちょうど$k$個あれば,一般項はそれらの冪の組み合わせになる.
  \item 根が重なれば固有ベクトルが足りず,\sectionref*{節：差分で線形漸化式を解く}のように倍率を取り除いて差分を追う.
\end{itemize}
どれだけ項数が増えても,行列を公比のように見て,掛けても定数倍になる向きを探す,という考え方は変わらないのである.

\subsection{角度を足す更新と,角度を倍にする更新}\label{節：角度の整数倍と多項式}

\sectionref*{節：非実数の特性根と回転}の
\[
  a_{n+2}=2\cos\theta\,a_{n+1}-a_n
\]
は,\,$a_n=\cos(n\theta)$のように\textbf{一定の角を足していく}更新であった.これは線形の漸化式である.
一方,本編\sectionref{節：三角関数型}の$a_{n+1}=2a_n^2-1$は,\,$a_n=\cos\theta_n$と置くと$\theta_{n+1}=2\theta_n$となる,\textbf{角度を倍にする}更新であった.こちらは線形ではない.
同じ余弦が現れるが,全く別の操作である.
この2つを,「整数倍の角を表す多項式」でつないでみよう.

$x=\cos\theta$とすると,\bookterm{addition}{加法定理}から
\[
  \cos((n+1)\theta)=2x\cos(n\theta)-\cos((n-1)\theta).
\]
同じ規則で多項式の列を
\begin{align*}
  &T_0(x)=1,\qquad T_1(x)=x,\\
  &T_{n+1}(x)=2xT_n(x)-T_{n-1}(x)
  \qquad(n\ge1)
\end{align*}

と定める.
最初から足し算と掛け算だけで作るので,すべて多項式になる.これを\bookterm{chebyshev}{チェビシェフ多項式}という.

この定義から
\[
  T_2(x)=2x^2-1,\qquad
  T_3(x)=4x^3-3x
\]
が得られる.
二倍角・三倍角の公式と同じ形である.一般にも,\,$n=0,1$の値と更新規則が一致するので,帰納法から
\[
  T_n(\cos\theta)=\cos(n\theta).
\]
整数倍の角の余弦を,\,$x=\cos\theta$の多項式で表せた.図は$T_2$による二倍角の横座標を示す.
\BookFigChebyshev

ここで本編の$a_1=\frac{\sqrt3}{2}$, $a_{n+1}=2a_n^2-1$に戻ろう.
本編では$a_1=\cos(\frac{\pi}{6})$から角度を二倍ずつ進めたが,多項式の言葉では
\[
  a_{n+1}=T_2(a_n),\qquad
  a_n=T_{2^{n-1}}(a_1)
\]
と書ける.
初項そのものの多項式で\bookterm{general}{一般項}を表せた.その理由は次の合成の関係にある.

一般に,\,$x=\cos\theta$なら
\BookMathLayout{\[
    T_m(T_k(x))
    =T_m(\cos(k\theta))
    =\cos(mk\theta)
    =T_{mk}(x).
  \]}{
  \[
    \begin{aligned}
      T_m(T_k(x))&=T_m(\cos(k\theta))\\
      &=\cos(mk\theta)=T_{mk}(x).
    \end{aligned}
  \]
}
この等式は $[-1,1]$ のすべての $x$ で成り立ち,両辺が多項式なので,多項式の恒等式として実数全体で成り立つ.
したがって$T_2$を繰り返し合成する操作は,添字を二倍ずつ進める$T_2,T_4,T_8,\ldots$にまとまる.
本編の「整数倍の公式へ」で扱った$m$倍角にも同じことがいえる.
整数$m\ge2$に対して$a_{n+1}=T_m(a_n)$なら
\[
  a_n=T_{m^{n-1}}(a_1)
\]
である.
$-1\le a_1\le1$なら$a_1=\cos\theta$と表せるため,これは本編の$a_n=\cos(m^{n-1}\theta)$と一致する.
多項式としての式は,\bookterm{trigonometric}{三角関数}で表せない初項についても成り立つ.

\paragraphbf{三角関数で表せない初項：双曲線関数}\label{節：双曲線関数とチェビシェフ多項式}
では,\,$|a_1|>1$のとき,\,$a_n=T_{m^{n-1}}(a_1)$は何を表しているのだろうか.
ここで,次の関数を導入しよう.
\[
  \cosh x=\frac{e^x+e^{-x}}{2}
\]
これを双曲線余弦関数という.\footnote{$\sinh x=\frac{e^x-e^{-x}}{2}$と合わせて$\cosh^2x-\sinh^2x=1$を満たす.点$(\cos x,\sin x)$が円$X^2+Y^2=1$を描くのに対し,点$(\cosh x,\sinh x)$は双曲線$X^2-Y^2=1$を描くので,この名前がある.}
$\cosh x\ge1$であり,1以上のどんな値も$\cosh x$の形で表せる.
定義から直接計算すると,
\begin{align*}
  &\cosh((n+1)x)+\cosh((n-1)x)\\
  &\qquad=2\cosh x\cosh(nx)
\end{align*}
が成り立つ.余弦の加法定理から得た式と全く同じ形である.
したがって,\,$x=\cosh t$と置いても$T_n$と同じ規則で値が進み,
\[
  T_n(\cosh t)=\cosh(nt)
\]
となる.特に$\cosh2t=2\cosh^2t-1$は,倍角の公式そのものである.
$a_1>1$なら$a_1=\cosh t$と書けるので,\,$a_{n+1}=T_m(a_n)$の一般項は$a_n=\cosh(m^{n-1}t)$になるのだ.
$\cos$が角度を整数倍するのに対し,\,$\cosh$は$t$を整数倍する.

この見方で,本編の\sectionref*{節：対称式を使う}を振り返ろう.
\[
  a_1=3,\qquad a_{n+1}=a_n^2-2
\]
は,\,$\frac{a_n}{2}$について$\frac{a_{n+1}}{2}=2\left(\frac{a_n}{2}\right)^2-1=T_2\left(\frac{a_n}{2}\right)$を満たす.
つまり,三角関数型の$a_n$を2倍しただけの漸化式だったのである.
$\frac{a_1}{2}=\frac32>1$なので,\,$a_n=2\cosh t_n$と置けば$t_{n+1}=2t_n$であり,
\[
  a_n=2\cosh\left(2^{n-1}t_1\right),\qquad e^{t_1}=\frac{3+\sqrt5}{2}
\]
を得る.\footnote{$2\cosh t_1=3$は$e^{t_1}+e^{-t_1}=3$なので,\,$e^{t_1}$は$u^2-3u+1=0$の解である.}
$2\cosh t=e^t+e^{-t}$だから,\,$u=e^t$と置けばこれは$u+\frac1u$にほかならない.
本編で「積が1になる2数の和」という置き換えを使ったのは,\,$\cosh$の倍角の公式を使っていたことにあたるのだ.
一般に
\[
  2T_m\left(\frac{u+u^{-1}}{2}\right)=u^m+u^{-m}
\]
が成り立つので,\,$a_{n+1}=a_n^3-3a_n$のような3倍角の漸化式も同じく解ける.\footnote{$m=3$では$2T_3\left(\frac a2\right)=a^3-3a$である.}

実数の初項で整理しておこう.
\begin{itemize}
  \item $-2\le a_1\le2$なら$a_1=2\cos\theta$と書け,\,$a_n=2\cos(2^{n-1}\theta)$となる.
  \item $a_1\ge2$なら$a_1=2\cosh t$と書け,\,$a_n=2\cosh(2^{n-1}t)$となる.
  \item $a_1<-2$なら$a_1=-2\cosh t$と書ける.\,$a_2=a_1^2-2=2\cosh2t$で符号が消えるので,\,$n\ge2$では$a_n=2\cosh(2^{n-1}t)$である.
\end{itemize}
三角関数型と基本対称式型は,初項が$[-2,2]$の内側か外側かで呼び名が変わっただけの,同じ漸化式なのである.

\paragraphbf{小話：リュカ数}
$\frac{3+\sqrt5}{2}=\left(\frac{1+\sqrt5}{2}\right)^2$なので,\,$a_1=3$の場合の各項$3,\allowbreak 7,\allowbreak 47,\allowbreak 2207,\allowbreak \ldots$は,黄金比$\varphi=\frac{1+\sqrt5}{2}$を使って$a_n=\varphi^{2^n}+\varphi^{-2^n}$と書ける.
$L_k=\varphi^k+\left(-\frac1\varphi\right)^k$で定まる数列$2,\allowbreak 1,\allowbreak 3,\allowbreak 4,\allowbreak 7,\allowbreak 11,\allowbreak 18,\allowbreak 29,\allowbreak 47,\allowbreak \ldots$はリュカ数と呼ばれ,フィボナッチ数列と同じ漸化式$L_{k+2}=L_{k+1}+L_k$を満たす.
$a_n=L_{2^n}$であり,この漸化式はリュカ数の$2,4,8,16,\ldots$番目を順に取り出していたのだ.


多項式列を\bookterm{recurrence}{漸化式}で作る例は,\bookterm{chebyshev}{チェビシェフ多項式}だけではない.
\bookterm{legendre}{ルジャンドル多項式} $P_n(x)$ は,初期値
\[
  P_0(x)=1,\qquad P_1(x)=x
\]
と三項\bookterm{recurrence}{漸化式}
\BookMathLayout{
  \[
    \begin{aligned}
      (n+1)P_{n+1}(x)&=(2n+1)xP_n(x)\\
      &\quad-nP_{n-1}(x)\qquad(n\ge1)
    \end{aligned}
  \]
}{
  \[
    \begin{aligned}
      (n+1)P_{n+1}(x)&=(2n+1)xP_n(x)\\
      &\quad-nP_{n-1}(x)\qquad(n\ge1)
    \end{aligned}
  \]
}
で定まる.
例えば
\[
  P_2(x)=\frac{3x^2-1}{2}
\]
である.
\bookterm{chebyshev}{チェビシェフ多項式}と同じく三項\bookterm{recurrence}{漸化式}から作られるが,\bookterm{legendre}{ルジャンドル多項式}では係数が $n$ にも依存する.

\bookterm{legendre}{ルジャンドル多項式}は,区間 $[-1,1]$ 上で重み $1$ に関して直交する.
ここでいう直交性とは,異なる次数の二つを掛けて区間全体で積分すると $0$ になる性質である.
\[
  \int_{-1}^{1}P_m(x)P_n(x)\,dx=0
  \qquad(m\ne n)
\]
この性質は関数の表示や近似に役立つ.ここでは,三項\bookterm{recurrence}{漸化式}が生む別の多項式列として紹介する.

\subsection*{コラム：有限個の点から多項式を作る}

最初の三項が $1,2,3$ なら,
$x$ と $x+(x-1)(x-2)(x-3)$ はどちらもその三項に一致するが,
$n=4$ ではそれぞれ $4$ と $10$ になる.
有限個の値が一致しても,その後の値まで同じとは限らない.
与えられた値を通る多項式を,一般の場合にも作ってみよう.

最初の$N$項$a_1,\ldots,a_N$に対して,
\[
  L(x)=\sum_{k=1}^{N}a_k
  \prod_{\substack{1\le j\le N\\j\ne k}}
  \frac{x-j}{k-j}
\]
と置くと, $L(k)=a_k$ $(1\le k\le N)$ となる.
$x=k$ のとき $k$ 番目の積は $1$ になり,ほかの積には $x-k$ の因子があるためである.
次数を $N-1$ 以下に制限した多項式の中では,この $L$ は一意である.二つの候補の差は $1,\ldots,N$ を根にもつ次数 $N-1$ 以下の多項式なので,零多項式になる.

次数制限はデータからは決まらない.\,$L(x)+c\prod_{j=1}^{N}(x-j)$も,任意の定数$c$で最初の$N$個の値を保つ.
次の図は,別の三点を通り,途中の形が異なる二つの曲線を描いた例である.
\BookFigLagrange

\bookterm{interpolation}{ラグランジュ補間}は一つの延長候補を作る.それを元の規則と結論づけるには,次数や\bookterm{recurrence}{漸化式}など追加の条件が要る.
